\documentclass[tikz,border=8pt]{standalone}
\usepackage{tikz}
\usepackage{amsmath}
\usepackage{amssymb}
\usepackage{pgfplots}
\pgfplotsset{compat=1.18}
\usepackage{ctex}
\usetikzlibrary{calc,positioning,arrows.meta,trees}

% LC 104 最大深度（Maximum Depth of Binary Tree）
% 树：[3,9,20,null,null,15,7]
% 每节点旁注释递归返回值
% 叶子返回 1，内部节点返回 max(left,right)+1

\begin{document}
\begin{tikzpicture}[
  every node/.style={},
  level distance=1.6cm,
  level 1/.style={sibling distance=4.2cm},
  level 2/.style={sibling distance=2.0cm},
  edge from parent/.style={draw,->,>=Stealth,thick}
]

%% ── 标题 ──────────────────────────────────────────────────
\node[font=\large\bfseries] at (0, 1.9) {LC 104 二叉树最大深度（DFS 后序）};
\node[font=\small,gray] at (0, 1.35) {递归：\texttt{maxDepth(node) = max(left,right)+1}，叶子 = 1};

%% ── 主树 ─────────────────────────────────────────────────
\node[draw,circle,minimum size=0.78cm,fill=blue!10,font=\bfseries] (root) at (0,0) {3}
  child {
    node[draw,circle,minimum size=0.78cm,fill=blue!10,font=\bfseries] (n9) {9}
  }
  child {
    node[draw,circle,minimum size=0.78cm,fill=blue!10,font=\bfseries] (n20) {20}
    child {
      node[draw,circle,minimum size=0.78cm,fill=blue!10,font=\bfseries] (n15) {15}
    }
    child {
      node[draw,circle,minimum size=0.78cm,fill=blue!10,font=\bfseries] (n7) {7}
    }
  };

%% ── 每节点返回值标注（右侧小方框） ──────────────────────
% 叶子 9：返回 1
\node[draw=gray!60,fill=orange!8,rounded corners=2pt,inner sep=3pt,
      font=\small,align=center]
  at (-4.2,-1.6) {返回 \textbf{1}\\（叶子）};
\draw[->,>=Stealth,gray!60,thin] (-3.3,-1.6) -- (-2.0,-1.6);

% 叶子 15：返回 1
\node[draw=gray!60,fill=orange!8,rounded corners=2pt,inner sep=3pt,
      font=\small,align=center]
  at (-0.5,-3.6) {返回 \textbf{1}};
\draw[->,>=Stealth,gray!60,thin] (0.4,-3.6) -- (1.0,-3.6);

% 叶子 7：返回 1
\node[draw=gray!60,fill=orange!8,rounded corners=2pt,inner sep=3pt,
      font=\small,align=center]
  at (3.8,-3.6) {返回 \textbf{1}};
\draw[->,>=Stealth,gray!60,thin] (2.9,-3.6) -- (2.3,-3.6);

% 节点 20：max(1,1)+1 = 2
\node[draw=gray!60,fill=green!8,rounded corners=2pt,inner sep=3pt,
      font=\small,align=center]
  at (4.8,-1.6) {$\max(1,1)+1$\\$= \mathbf{2}$};
\draw[->,>=Stealth,gray!60,thin] (3.6,-1.6) -- (2.5,-1.6);

% 根 3：max(1,2)+1 = 3
\node[draw=gray!60,fill=blue!8,rounded corners=2pt,inner sep=4pt,
      font=\small,align=center]
  at (4.2, 0.0) {$\max(1,2)+1$\\$= \mathbf{3}$};
\draw[->,>=Stealth,gray!60,thin] (2.9, 0.0) -- (2.0, 0.0);

%% ── null 节点标注 ─────────────────────────────────────────
\node[font=\scriptsize,gray] at (-2.1,-3.6) {\texttt{null}};
\node[draw,dashed,circle,minimum size=0.5cm,gray!40] at (-2.1,-3.4) {};
\node[font=\scriptsize,gray] at (0.2,-4.1) {$\to 0$};

\node[font=\scriptsize,gray] at (2.1,-3.6) {\texttt{null}};
\node[draw,dashed,circle,minimum size=0.5cm,gray!40] at (2.1,-3.4) {};
\node[font=\scriptsize,gray] at (3.8,-4.1) {$\to 0$};

%% ── 递归调用栈示意 ────────────────────────────────────────
\node[draw=gray!50,fill=white,thin,rounded corners=3pt,
      font=\small,inner sep=6pt,align=left]
  at (0,-5.6)
  {\textbf{后序递归回溯过程}：\\
   \quad $\texttt{maxDepth}(9) \;=\; 1$（叶子）\\
   \quad $\texttt{maxDepth}(15) = 1$（叶子）\\
   \quad $\texttt{maxDepth}(7) \;= 1$（叶子）\\
   \quad $\texttt{maxDepth}(20) = \max(1,1)+1 = 2$\\
   \quad $\texttt{maxDepth}(3) \;= \max(1,2)+1 = \mathbf{3}$};

\end{tikzpicture}
\end{document}
